\section{Discussion}\label{sec:discussion}
The central result is the coexistence of exact local mathematics
and negative complete-operation timing evidence. The original-row
updates preserve selected formulations without explicit departure
variables in the principal system. That algebraic property is
verified by independent equivalent formulations and original-space
checks. The registered experiments nevertheless fail to establish
their proposed timing gains. There is no contradiction: a smaller
principal representation is only one contributor to preparation,
iteration, recovery and validation cost.

The grouped operation adds a sorted threshold and joint dual
checks. A sparse reference may benefit from simple diagonal
auxiliaries or effective symbolic ordering. A prepared active-set
backend may exploit vector updates efficiently. Dense/sparse
conversion and language-level overhead can also matter. Without
separate controlled instrumentation, these are plausible mechanisms
to investigate rather than measured causal decompositions of the
observed difference. Numerical optimization and floating-point
analysis caution against inferring cost or accuracy from dimension
alone~\cite{nocedal,higham}.

Including the strongest compatible formulations is therefore
important. A symmetric grouped projection backend after directional
row splitting can solve the same directional objective. An epigraph
can represent the same group price with one auxiliary per group.
The computational comparison would be weaker if it omitted either
because its representation differs from the candidate. Similarly,
DAQP's existing side weights prevent an unequal-cost example from
supporting a new expressibility claim. Prior orthant and weighted
projection work constrains the interpretation of the analytic
construction as well as its implementation~\cite{orthant,weighted}.

Operational interpretation requires another separation. A finite
price can intentionally purchase an interval departure. Numerical
optimality does not prove that the chosen price reflects a real
engineering risk or that the resulting physical policy is safe.
Predictive control relies on model, constraints and closed-loop
assumptions beyond solving one convex subproblem~\cite{mayne,rawlings}.
The source-faithful and authored building examples are synthetic,
and the repeated forecasts are not measured building savings.
Their value here is transparent interpretation and repeatable
numerical comparison, not a field-performance claim.

The negative results constrain subsequent research. A revised
method should identify a genuinely different computational
mechanism, state why the strongest reference might fail to
exploit it, and fix a prospective population before measurement.
Selecting favorable current instances, weakening the oracle,
removing acceptance work or changing the failed-pair rule would
answer another question. The complete retained outcomes instead
make it possible to formulate that next question without
misrepresenting either study's original evidence.

Within the stated domain, the proximal formulas remain useful
tools for modeling and independent verification. Their explicit
zero-price, endpoint, sign and metric conditions clarify when
an implementation is solving the intended problem. The present
evidence supports those restricted mathematical and computational
statements. It supplies no global priority claim, general runtime
superiority or guarantee of practical benefit outside the tested
synthetic regimes.
